This research project focusses on software, namely the development of more efficient data-intensive processing pipelines applicable in machine learning and neural network training. The Green Software Foundation \footnote{Foundation's website available at \url{https://greensoftware.foundation/}} states that software may be "part of the climate problem or part of the climate solution." Green Software is described as software that minimises its environmental effect by using fewer physical resources and less energy, while optimising energy use to employ lower-carbon sources \cite{WhatGreenSoftware2021}. In the realm of \gls{ML} and \gls{NN} training, minimising training duration—and therefore the read and write latency associated with the dataset—has been shown to significantly decrease carbon emissions~\cite{pattersonCarbonEmissionsLarge2021,pattersonCarbonFootprintMachine2022}. This thesis minime latency and thus enhance data throughput for reading and writing on Iceberg tables on \gls{HopsFS}. This objective adheres to the essential principles of green software by reducing \gls{CPU} time utilisation relative to the prior system. Minimising \gls{CPU} utilisation lowers energy consumption, resulting in a reduced carbon impact. This as a concequence, enhance energy efficiency of data-intensive computer pipelines, which are extensively used in \gls{ML} and \gls{NN} training.

For those reasons, this project also supports the \glspl{SDG}~\footnote{\glspl{SDG} website available at \url{https://sdgs.un.org/}} Affordable and Clean Energy (Goal 7) and Industry Innovation and Infrastructure (Goal 9), particularly the objectives 7.3, "Double the improvement in energy efficiency", and 9.4, "Upgrade all industries and infrastructures for sustainability".

\begin{figure}
    \centering
    \begin{subfigure}[b]{0.5\linewidth}
        \centering
        \includegraphics[width=0.6\linewidth]{figures/1-introduction/SDG_icons-07.png} 
        \label{fig:sdg07}
    \end{subfigure}\hfill
    \begin{subfigure}[b]{0.5\linewidth}
        \centering
        \includegraphics[width=0.6\linewidth]{figures/1-introduction/SDG_icons-09.png}
        \label{fig:sdg09}
	\end{subfigure}
	\caption[Sustainable Development Goals supported by this thesis]{Illustrations of the \glstext{SDG} supported by this thesis.}
	\label{fig:sdgs}
\end{figure}

Furthermore, the experiments created to satisfy \gls{G}3 has been carefully design to minimize the number of trails necessary for statistical relevance, pursuing resource effiency. All the data throughly explained and all the codes are and provided in this thesis report, to best address reproducibility of this project.